%% bare_jrnl.tex
%% V1.4b
%% 2015/08/26
%% by Michael Shell
%% see http://www.michaelshell.org/
%% for current contact information.
%%
%% This is a skeleton file demonstrating the use of IEEEtran.cls
%% (requires IEEEtran.cls version 1.8b or later) with an IEEE
%% journal paper.
%%
%% Support sites:
%% http://www.michaelshell.org/tex/ieeetran/
%% http://www.ctan.org/pkg/ieeetran
%% and
%% http://www.ieee.org/
%%
%% ---------------------------------------------------------------
%% Modified by PaperCrew (TFG).
%% Plantilla de referencia: IEEE Journal LaTeX Template oficial
%% (IEEEtran.cls, distribuido vía template-selector.ieee.org),
%% con el front-matter ajustado a las Author Guidelines de
%% IEEE Internet of Things Journal (ieee-iotj.org/guidelines-for-authors).
%% Abstract e introducción se insertan automáticamente.
%% ---------------------------------------------------------------

%%*************************************************************************
%% Legal Notice:
%% This code is offered as-is without any warranty either expressed or
%% implied; without even the implied warranty of MERCHANTABILITY or
%% FITNESS FOR A PARTICULAR PURPOSE!
%% User assumes all risk.
%% In no event shall the IEEE or any contributor to this code be liable for
%% any damages or losses, including, but not limited to, incidental,
%% consequential, or any other damages, resulting from the use or misuse
%% of any information contained here.
%%
%% All comments are the opinions of their respective authors and are not
%% necessarily endorsed by the IEEE.
%%
%% This work is distributed under the LaTeX Project Public License (LPPL)
%% ( http://www.latex-project.org/ ) version 1.3, and may be freely used,
%% distributed and modified. A copy of the LPPL, version 1.3, is included
%% in the base LaTeX documentation of all distributions of LaTeX released
%% 2003/12/01 or later.
%% Retain all contribution notices and credits.
%% ** Modified files should be clearly indicated as such, including  **
%% ** renaming them and changing author support contact information. **
%%*************************************************************************


\documentclass[journal]{IEEEtran}

\hyphenation{op-tical net-works semi-conduc-tor}


\begin{document}

\title{[ARTICLE TITLE]}

\author{[FIRST AUTHOR],~\IEEEmembership{Member,~IEEE,}
        and~[CO-AUTHOR]%
\thanks{[AFFILIATION DATA].}%
\thanks{Manuscript received [DATE].}}

\markboth{IEEE INTERNET OF THINGS JOURNAL,~Vol.~[VOL], No.~[NUM], [MONTH]~[YEAR]}%
{[AUTHORS]: [SHORT TITLE]}

\maketitle

% IEEE IoT-J: el abstract debe tener entre 150 y 250 palabras y no debe
% incluir citas bibliográficas ni ecuaciones.
\begin{abstract}
% [Pendiente de generacion]
\end{abstract}

\begin{IEEEkeywords}
[KEYWORD 1], [KEYWORD 2], [KEYWORD 3]
\end{IEEEkeywords}

\IEEEpeerreviewmaketitle


\section{Introduction}
\label{sec:intro}

El procesamiento del lenguaje natural (NLP) ha avanzado significativamente gracias a la incorporación de técnicas de pre-entrenamiento de representaciones del lenguaje basadas en grandes cantidades de texto no etiquetado. Este enfoque ha demostrado ser fundamental para mejorar el desempeño en una gran variedad de tareas lingüísticas, desde la clasificación y el análisis semántico hasta la inferencia y la respuesta a preguntas. Dentro de este panorama, se destaca la importancia de desarrollar modelos que puedan capturar eficazmente el contexto tanto a la izquierda como a la derecha de cada palabra o token, superando las limitaciones de los modelos unidireccionales tradicionales.

En la evolución del estado del arte, diversos métodos han explorado la construcción de representaciones profundas y contextualizadas del lenguaje, pero con restricciones en la manera en que se integra la información contextual o en la forma de adaptar las arquitecturas para tareas específicas. La bidireccionalidad en el modelado del contexto emerge como un aspecto crucial para representar de manera más rica y holística el significado y las relaciones lingüísticas, especialmente en tareas complejas como la inferencia de lenguaje natural y la comprensión de preguntas. Esto subraya la relevancia de avanzar hacia modelos de lenguaje que puedan incorporar esta bidireccionalidad de forma integrada y eficiente.

Los trabajos previos presentan limitaciones significativas que justifican la necesidad de la investigación propuesta. Peters et al. (2018) (ELMo)~\cite{peters2018elmo} ofrecen un modelo de representaciones contextuales bidireccionales basado en LSTM; sin embargo, su enfoque combina capas independientes en lugar de un modelo integral bidireccional, lo que limita la profundidad y la calidad de la información contextual integrada. Además, su método no proporciona un modelo unificado para múltiples tareas ni aprovecha el pre-entrenamiento de transformadores bidireccionales con objetivos de enmascaramiento, como propone BERT.

Por otro lado, Radford et al. (2018) [VERIFICAR] (OpenAI GPT) utilizan un modelo unidireccional de lenguaje basado en transformadores para el pre-entrenamiento generativo seguido de un ajuste fino supervisado. Sin embargo, esta unidireccionalidad restringe el contexto considerado para la representación de palabras, lo que resulta en una menor eficacia para tareas que requieren información contextual de ambos lados del token. Además, la arquitectura empleada por Radford et al. requiere adaptaciones específicas más complejas para cada tarea en comparación con el enfoque más generalizable de BERT.

Estas limitaciones en los enfoques previos, tanto en la falta de un modelado profundo y bidireccional integrado como en la necesidad de adaptaciones complejas por tarea, evidencian un hueco importante que el presente trabajo busca superar mediante el desarrollo de un modelo de Transformer bidireccional profundo con enmascaramiento y predicción de la siguiente oración, que permita una mejor representación contextual y transferencia generalizada a múltiples tareas.

La propuesta central del trabajo consiste en desarrollar BERT (Bidirectional Encoder Representations from Transformers), un modelo que aborda directamente el hueco identificado al pre-entrenar representaciones de lenguaje profundas y bidireccionales mediante dos tareas principales: el modelo de lenguaje enmascarado (Masked Language Model, MLM) y la predicción de la siguiente oración (Next Sentence Prediction, NSP). El uso del MLM posibilita que el modelo integre contexto tanto a la izquierda como a la derecha de cada token, superando la restricción presente en modelos unidireccionales, mediante el enmascaramiento aleatorio de tokens y la predicción de los tokens originales. Por su parte, la tarea de NSP ayuda a aprender relaciones entre pares de oraciones, lo cual es fundamental para tareas de inferencia y comprensión más complejas.

Este enfoque integral tiene sentido porque combina el poder de la bidireccionalidad completa con tareas de pre-entrenamiento que favorecen la modelación del lenguaje en profundidad y la comprensión de relaciones textuales, sin requerir modificaciones arquitectónicas específicas para cada tarea. Así, BERT logra representar el contexto de manera más rica y generalizable, permitiendo que se adapte eficazmente a diversas tareas de procesamiento de lenguaje natural mediante un proceso de afinamiento fino sencillo.

- La demostración de la importancia del pre-entrenamiento bidireccional en el contexto profundo utilizando un modelo Transformer.  
- Un modelo preentrenado que puede ser afinado con una única capa adicional para múltiples tareas de NLP de nivel oración y token, sin necesidad de arquitecturas específicas por tarea.  
- La introducción conjunta del objetivo de enmascaramiento y la tarea de predicción de la siguiente oración en el aprendizaje previo.  
- Resultados de vanguardia en once tareas de NLP, con mejoras absolutas significativas en benchmarks como GLUE, SQuAD y MultiNLI.  
- Evaluación detallada que muestra la superioridad del enfoque frente a modelos previos basados en LSTM o transformadores unidireccionales.

El trabajo se valida mediante un proceso de fine-tuning del modelo preentrenado en once tareas variadas de procesamiento de lenguaje natural, que abarcan desde inferencia textual y respuesta a preguntas hasta análisis semántico, reconocimiento de entidades nombradas y clasificación de texto. Esta amplia gama de experimentos proporciona un respaldo sólido a las contribuciones presentadas, permitiendo evaluar la eficacia y generalización del modelo en diferentes contextos y tipos de tareas. Además, se realizan estudios de ablación que examinan el impacto de componentes específicos del modelo y su tamaño, contribuyendo a comprender mejor las fuentes de mejora en el desempeño.

El resto del artículo se organiza en las siguientes secciones: introducción, descripción detallada del modelo BERT, revisión de trabajos relacionados, explicación de los procedimientos de pre-entrenamiento y fine-tuning, presentación de los experimentos y análisis de resultados, y finalmente, información suplementaria en los apéndices. Cada una de estas secciones desarrolla de manera específica distintos aspectos del método, el contexto teórico, la validación empírica y los detalles técnicos del trabajo.


\section{[SECTION 2 TITLE]}
\label{sec:section2}

[Section 2 content.]


\section{[SECTION 3 TITLE]}
\label{sec:section3}

[Section 3 content.]


\section{[SECTION 4 TITLE]}
\label{sec:section4}

[Section 4 content.]


\section{Conclusion}
\label{sec:conclusion}

[Conclusions.]


\appendices
\section{Appendix}
[Appendix content.]


\section*{Acknowledgment}

[Acknowledgments.]


\bibliographystyle{IEEEtran}
\bibliography{bert}


\begin{IEEEbiographynophoto}{[FIRST AUTHOR]}
[Brief author biography.]
\end{IEEEbiographynophoto}


\end{document}
